\paragraph{Multi-agent communication.}
Agent communication languages~\cite{finin_kqml_1994} defined formal message types, and learned communication~\cite{sukhbaatar_commnet_2016, foerster_dial_2016} later optimized bandwidth end-to-end. Recent theoretical work examines complementary dimensions: when communication helps~\cite{rizvi_martel_communication_2026}, what to communicate~\cite{chen_five_ws_2026}, and why agents communicate~\cite{zhou_why_agents_communicate_2025}. Multi-agent debate~\cite{du_debate_2024} and LLM orchestration frameworks~\cite{li_camel_2023, hong_metagpt_2024, wu_autogen_2024, chen_agentverse_2024, wang_megaagent_2025} confirm that communication structure affects downstream performance. These contributions operate at the system level; we work at the message level, classifying by communicative function and testing whether functional categories predict pruning safety.

\paragraph{Communication efficiency.}
Existing efficiency methods operate at granularities coarser than message function. Agent-level pruning~\cite{wang_agentdropout_2025, wang_agentdropout_v2_2026, zhang_agentprune_2025, li_agp_2025, xiao_agentdiet_2026} removes entire agents with substantial token reduction. Turn-level stopping~\cite{zeng_s2mad_2025} halts conversations early. Topology optimization~\cite{chen_optima_2025, chen_goagent_2026, wang_anymac_2025, li_lacp_2025} restructures agent organization. Token- and trajectory-level filtering~\cite{safesieve_2025, lin_supervisoragent_2026} selectively removes inter-agent content by combining safety and relevance filters (dual-mechanism) or LLM-free scoring. Latent-space communication~\cite{zou_latentmas_2025, du_interlat_2025} and holistic optimization~\cite{li_omac_2025, jiang_agentcomm_2026} pursue complementary directions (Appendix Table~\ref{tab:method_comparison}). The shared limitation across these approaches is that none link communicative function to pruning outcomes. Our taxonomy-driven pruning identifies which message categories are safe to suppress based on their functional role, not their surface form or position.

\paragraph{Theoretical foundations.}
Our taxonomy draws on speech act theory~\cite{austin_how_1962, searle_speech_acts_1969} and dialogue act schemes (DAMSL~\cite{core_damsl_1997}, ISO 24617-2~\cite{bunt_iso_2012}, Switchboard DA~\cite{stolcke_dialogue_acts_2000}). Backchannels~\cite{yngve_getting_1970}, continuers~\cite{schegloff_discourse_1982}, and acknowledgment tokens~\cite{jefferson_acknowledgment_1984} carry minimal propositional content yet serve a load-bearing role in conversational grounding~\cite{clark_contributing_1989}. Computational backchannel studies~\cite{inoue_backchannel_2025, qian_backchannel_2026} confirm that these signals persist in neural dialogue systems with measurable effects on turn coordination. Extensions to LLM settings show that standard dialogue act categories apply to LLM conversations~\cite{qamar_dialogue_acts_2025} and that acknowledgment serves key functional roles in dual-agent mathematical reasoning~\cite{zhang_communication_strategies_2025}. These studies describe communicative strategies but do not perform causal interventions to test whether functional categories predict pruning safety, nor do they validate across multiple backbone and domain configurations. Attribution methods~\cite{xia_hivemind_2026, horovicz_agentshap_2025, cemri_mast_2025} assign credit at the agent or trace level but do not decompose individual messages by communicative function. We extend these foundations to inter-agent communication and connect functional categories directly to pruning outcomes through causal intervention and multi-configuration statistical validation.
